% Inactive research fragment; not an active author input.
% Standard inputs: convexity and positive-matrix monotonicity of the
% spectral Garding cone, bounded psh Liouville, and interior Laplace estimates.
% This fragment does not prove a derivative-independent second-order estimate.
\section{Liouville and gradient estimates for the Hessian cone}

Let $1\leq m\leq n$, and let $C_{n,m}$ be the closed cone of Hermitian matrices whose eigenvalues belong to $\overline{\Gamma_m}$, with $\Gamma_m$ as in the preceding fragment.
We use the standard algebraic facts that this spectral cone is convex, contains the nonnegative matrices, and is preserved by adding such matrices. Adding $\delta I$, $\delta>0$, to a matrix in $C_{n,m}$ gives a matrix in its interior. The polynomial $S_m$ is positive and increases under addition of a nonnegative matrix in that interior. These are algebraic inputs; no bounded Hessian-product theorem is assumed below.

Write $H(v)=(v_{i\bar j})$. A continuous real function $v$ on a Euclidean domain is called admissible here if
\[
 \operatorname{tr}(QH(v))\geq0
 \quad\hbox{in distributions for every }Q\in C_{n,m}^*,
\]
where $C_{n,m}^*=\{Q:\operatorname{tr}(QA)\geq0\text{ for all }A\in C_{n,m}\}$ is the dual cone.
For smooth functions this says precisely that $H(v)\in C_{n,m}$.
We call an admissible function maximal for this cone if every smooth real function $h$ touching it from below satisfies $H(h)\notin\operatorname{int}C_{n,m}$ at the contact point.
Admissibility implies subharmonicity, because $\operatorname{tr}A\geq0$ on $C_{n,m}$. It is preserved by convolution with a nonnegative smooth kernel. For smooth admissible functions, nonnegative weighted averages of their Hessians also belong to the cone.

\begin{lemma}\label{draft:596-hessian-cone-stability}
Local uniform limits of admissible maximal functions are admissible and maximal. Suppose $v$ is admissible and maximal on a bounded open set $D$, and $h$ is smooth and admissible there. If both are continuous on $\overline D$ and $h\leq v$ on $\partial D$, then $h\leq v$ on $D$.
\end{lemma}
\begin{proof}
The distributional inequalities pass to local uniform limits. If a lower test $h$ for the limit at $z_0$ has Hessian in the interior, subtract $\delta|z-z_0|^2$, for sufficiently small $\delta>0$, to make the contact strict while retaining interior admissibility near $z_0$. Interior minimum points for the approximants then tend to $z_0$. Adding the minimum value produces lower tests for them with Hessian in the interior, contradicting their maximality.

For comparison choose $R$ with $\overline D\subset\{|z|\leq R\}$. If $h>v$ somewhere, the function $h_\delta=h+\delta(|z|^2-R^2)$ still exceeds $v$ somewhere for small $\delta>0$ and is at most $v$ on the boundary. The negative minimum of $v-h_\delta$ is attained inside $D$. Since $H(h_\delta)=H(h)+\delta I$ belongs to the interior, adding the minimum value gives a forbidden lower test. This argument requires no regularity of the boundary.
\end{proof}

\begin{theorem}\label{draft:596-hessian-cone-liouville}
Every bounded globally Lipschitz admissible function on $\mathbb C^n$ that is maximal for $C_{n,m}$ is constant.
\end{theorem}
\begin{proof}
For $m=1$, the distributional subharmonicity and the lower-test condition say that the function is a viscosity solution of the Laplace equation. The elementary equivalence with harmonicity and bounded harmonic Liouville give the conclusion. For $m=n$, admissibility is distributional plurisubharmonicity. The continuous representative is plurisubharmonic, and bounded plurisubharmonic Liouville applies.
For a fixed $m\geq2$, we therefore argue by induction on $n\geq m$.

The reduction in dimension is particularly simple for this cone. For a Hermitian $(n-1)$-by-$(n-1)$ matrix $A$ and $m<n$,
\[
 S_j(A\oplus0)=S_j(A),\qquad 1\leq j\leq m.
\]
Consequently $A\oplus0$ belongs to $C_{n,m}$, respectively its interior, exactly when $A$ belongs to $C_{n-1,m}$, respectively its interior. For the closed cones, the nonnegative inequalities characterize the closure: the expansion of $S_j(\lambda+\delta\mathbf1)$ has nonnegative coefficients and a strictly positive $\delta^j$ term. If an admissible function is independent of $z_n$, convolving it and using the block identity shows that the function in the remaining variables is admissible. Extending its lower tests independently of $z_n$ shows that maximality also passes to dimension $n-1$.

Suppose $v$ is nonconstant. Normalize $\inf v=0$, $\sup v=1$, and let $L$ be its Lipschitz constant. Choose a smooth radial nonnegative probability kernel $\eta$, supported in the closed unit ball and positive in its interior. Put
\[
 [f]_r(x)=\int f(x+y)\eta_r(y)\,\mathrm dy,
 \qquad \eta_r(y)=r^{-2n}\eta(y/r),\qquad v_\epsilon=[v]_\epsilon.
\]
Both $v$ and $v^2$ are subharmonic: for the second assertion use
$\Delta(v_\epsilon^2)=2v_\epsilon\Delta v_\epsilon+2|\nabla v_\epsilon|^2\geq0$ and pass to the uniform limit.
Their radial averages increase with $r$, by monotonicity of spherical means. For either $f=v$ or $f=v^2$,
\[
 |[f]_r(x)-[f]_r(y)|
 \leq\|f\|_\infty\|\eta_r(\cdot-x)-\eta_r(\cdot-y)\|_{L^1}
 \leq C\frac{|x-y|}{r}.
\]
Thus the limiting average is independent of its center. The submean inequality makes it at least $f(x)$ for every $x$, and the upper bound is $\sup f=1$. Hence
\[
 [v]_r(x)\longrightarrow1,\qquad [v^2]_r(x)\longrightarrow1.
\]

First suppose there are a fixed $\rho>0$, numbers $\epsilon_k\downarrow0$, $r_k\to\infty$, points $x_k$, and unit complex directions $\xi_k$, such that
\[
 [v^2]_{r_k}(x_k)+[v]_\rho(x_k)-2v(x_k)\geq\frac43,
 \qquad
 \int_{B(x_k,r_k)}|\partial_{\xi_k}v_{\epsilon_k}|^2\,\mathrm dy\longrightarrow0.
\]
Translate and rotate to make $x_k$ the origin and $\xi_k$ the last complex direction. Compactness of uniformly bounded Lipschitz functions and the preceding lemma give a locally uniform limit $a$ that is admissible and maximal. The transformed mollifications have the same limit. Their last derivatives tend to zero in $L^2$ on every fixed ball, so $\partial_{z_n}a=0$ in distributions. Since $a$ is real and continuous, it is independent of $z_n$.
The induction hypothesis gives $a=c\in[0,1]$. In the displayed gap inequality the first average is at most one, while the fixed-radius average and the value at the origin tend to $c$. This would give $4/3\leq1-c\leq1$, a contradiction.

If this case does not occur, then for every fixed $\rho>0$ there is $C_\rho\geq1$ such that
\[
 \int_{B(x,r)}|\partial_\xi v_\epsilon|^2\,\mathrm dy\geq C_\rho^{-1}
\]
for every unit complex direction $\xi$, whenever $0<\epsilon<C_\rho^{-1}$, $r>C_\rho$ and
$[v^2]_r(x)+[v]_\rho(x)-2v(x)\geq4/3$.
Indeed, failure for every integer choice $C_\rho=k$ produces the sequences excluded above. The dependence on the fixed radius $\rho$ is retained.
Translate so that $v(0)<1/12$. Choose $\rho$ with $[v]_\rho(0)>3/4$, then $r>C_\rho$ with $[v^2]_r(0)>3/4$. The open set
\[
 D=\{x:2v(x)<[v^2]_r(x)+[v]_\rho(x)-4/3\}
\]
contains the origin.

On this set, $[v^2]_r-c|z|^2$ is smooth and admissible for some fixed $c>0$. To prove this, decompose
\[
 H([v_\epsilon^2]_r)
 = [2v_\epsilon H(v_\epsilon)]_r
   +[2(\partial_i v_\epsilon\overline{\partial_jv_\epsilon})_{ij}]_r
 =A_\epsilon+B_\epsilon.
\]
The matrix $A_\epsilon$ belongs to $C_{n,m}$ by convexity and $v_\epsilon\geq0$.
For every complex direction the unweighted energy on $B(x,r)$ is at least $C_\rho^{-1}$. Since $|\nabla v_\epsilon|\leq L$, a sufficiently thin boundary annulus carries at most half that amount, uniformly in $x,\epsilon,\xi$. On the remaining smaller closed ball $\eta_r$ has a positive minimum. Therefore $B_\epsilon\geq c_0I$ on $D$, with $c_0>0$ independent of small $\epsilon$.
Choose $0<c\leq c_0$. Then $A_\epsilon+B_\epsilon-cI\in C_{n,m}$. Uniform convergence of $v_\epsilon^2$ and convolution at the fixed radius give convergence of every derivative to those of $[v^2]_r$, proving the assertion.

The smooth function
\[
 h=\tfrac12\bigl([v^2]_r-c|z|^2+[v]_\rho-4/3\bigr)
\]
is admissible on $D$. The set $D'=\{v<h\}$ contains the origin, is bounded, and lies in $D$. Continuity gives $v=h$ on its boundary. The comparison lemma gives $h\leq v$ on $D'$, a contradiction. This completes the induction.
\end{proof}

\begin{proposition}\label{draft:596-hessian-trace-to-gradient}
Fix a compact complex manifold $X$ of dimension $n$, a smooth Hermitian form $\omega$, a smooth real $(1,1)$-form $\theta$, an integer $2\leq m\leq n$, and constants $M,A,B>0$.
There is a uniform bound for $\sup_X|\partial u|_\omega$ over all smooth real $u$ such that $\theta_u=\theta+\ddc u\in\Gamma_m(\omega)$ and
\[
 |u|\leq M,\qquad
 \frac{\theta_u^m\wedge\omega^{n-m}}{\omega^n}\leq B,
 \qquad
 \operatorname{tr}_\omega\theta_u
       \leq A\bigl(1+\sup_X|\partial u|_\omega^2\bigr).
\]
The bound requires no derivatives of the Hessian density.
\end{proposition}
\begin{proof}
In local holomorphic coordinates write $\omega=\mathrm i g_{i\bar j}\,\mathrm dz_i\wedge\mathrm d\bar z_j$ and $\theta_u=\mathrm i(b_{i\bar j}+\kappa u_{i\bar j})\,\mathrm dz_i\wedge\mathrm d\bar z_j$, where the fixed constant $\kappa>0$ accounts for the convention for $\ddc$.
For eigenvalues $\lambda\in\Gamma_m$, $m\geq2$, one has
\[
 \sum_i\lambda_i^2
  =\left(\sum_i\lambda_i\right)^2-2S_2(\lambda)
  <\left(\sum_i\lambda_i\right)^2.
\]
Thus the assumed trace estimate bounds every eigenvalue in absolute value by $A(1+\sup|\partial u|^2)$.

Suppose contrary to the claim that a sequence $u_j$ has maxima $R_j=\sup|\partial u_j|_\omega\to\infty$, attained at $x_j$. Pass to $x_j\to x_\infty$ and work in a fixed chart, changing coordinates linearly so that $g(x_\infty)=I$. If $y_j$ is the coordinate of $x_j$, set
\[
 v_j(z)=u_j(y_j+z/R_j),\quad
 g_j(z)=g(y_j+z/R_j),\quad
 E_j(z)=R_j^{-2}b(y_j+z/R_j)+\kappa H(v_j)(z).
\]
These are defined on balls with radii tending to infinity. On each fixed ball $g_j\to I$, the smooth background term tends to zero, and $E_j$ is uniformly bounded by the eigenvalue estimate. In particular $H(v_j)$ and its Euclidean trace are uniformly bounded.
Interior estimates for the fixed Euclidean Laplacian give local uniform $W^{2,q}$ bounds for every finite $q$. A diagonal subsequence converges in $C^1$ on compact subsets to a bounded globally Lipschitz $v$, with $|\partial v(0)|=1$.

The Hermitian matrices $g_j^{-1/2}E_jg_j^{-1/2}$ belong to the interior of $C_{n,m}$ and differ from $\kappa H(v_j)$ by a term tending uniformly to zero on fixed balls. Pairing with any fixed matrix in the dual cone and passing in distributions proves that $v$ is admissible.
Its equation is
\[
 \frac{S_m(g_j^{-1/2}E_jg_j^{-1/2})}{\binom nm}
   =R_j^{-2m}\frac{\theta_{u_j}^m\wedge\omega^{n-m}}{\omega^n}
      (y_j+z/R_j)\leq BR_j^{-2m}.
\]
We show maximality by lower tests, rather than by an unproved convergence assertion for variable-background Hessian measures. If a smooth lower test $\psi$ at $z_0$ had $H(\psi)(z_0)$ in the interior, choose a small $\delta>0$ so that $H(\psi)-\delta I$ remains there near $z_0$. Strictifying the contact by subtracting $\delta|z-z_0|^2$ gives interior minimum points $z_j\to z_0$ for $v_j$ minus this test. At those points
\[
 E_j\geq R_j^{-2}b+\kappa(H(\psi)-\delta I).
\]
The matrix on the right, measured with $g_j$, is uniformly inside the cone. Monotonicity of $S_m$ under nonnegative addition gives a fixed positive lower bound for the left side of the preceding equation. This contradicts $BR_j^{-2m}\to0$.
Hence $v$ is maximal. The Liouville theorem makes it constant, contradicting $|\partial v(0)|=1$.
\end{proof}

Only the given trace estimate, the density upper bound and fixed smooth backgrounds are used in the rescaling. On a fixed rescaled ball, variation of the metric is of order $R_j^{-1}$, and the background form enters with factor $R_j^{-2}$. No local potential for the Hermitian metric is used. Obtaining the assumed trace estimate without derivatives of a prescribed approximation remains a separate obligation.
